\section{\texorpdfstring{The \(I\)-adic metric and completeness}{The I-adic metric and completeness}}

Let \(R\) be a ring, let \(V\) be an \(R\)-module, and let \(I\) be an ideal of \(R\).
Assume that
\[
    \bigcap_{i=1}^{\infty} I^i V = 0.
\]
We set \(I^0V:=V\).

Fix a real number \(\gamma\) such that
\[
    0<\gamma<1.
\]
For \(v\in V\), define
\[
    \|0\|_I=0.
\]
If \(v\neq 0\), then there exists a unique integer \(i\geq 0\) such that
\[
    v\in I^iV\setminus I^{i+1}V.
\]
We define
\[
    \|v\|_I=\gamma^i.
\]

Equivalently, if
\[
    \operatorname{ord}_I(v)
    :=
    \max\{\,i\geq 0\mid v\in I^iV\,\},
\]
then
\[
    \|v\|_I=\gamma^{\operatorname{ord}_I(v)}
    \qquad (v\neq 0).
\]

\begin{definition}[\(I\)-adic metric]
    Define
    \[
        d_I(v,w)=\|v-w\|_I,
        \qquad v,w\in V.
    \]
    We call \(d_I\) the \(I\)-adic metric on \(V\).
\end{definition}

\begin{proposition}
    The function \(d_I\) is a metric on \(V\). More precisely, for all
    \(u,v,w\in V\), one has
    \[
        d_I(v,u)
        \leq
        \max\{\,d_I(v,w),d_I(w,u)\,\}.
    \]
    In particular, \(d_I\) is a non-archimedean metric.
\end{proposition}

\begin{proof}
    It is clear that
    \[
        d_I(v,w)\geq 0
    \]
    and that
    \[
        d_I(v,w)=d_I(w,v).
    \]
    Moreover,
    \[
        d_I(v,w)=0
    \]
    if and only if \(v-w\in \bigcap_{i\geq 1} I^iV\). By assumption,
    \[
        \bigcap_{i\geq 1} I^iV=0,
    \]
    so \(d_I(v,w)=0\) if and only if \(v=w\).

    It remains to prove the triangle inequality. Let
    \[
        a=v-w,\qquad b=w-u.
    \]
    Then
    \[
        v-u=a+b.
    \]
    Suppose
    \[
        a\in I^rV,\qquad b\in I^sV.
    \]
    If \(m=\min\{r,s\}\), then
    \[
        a,b\in I^mV,
    \]
    hence
    \[
        a+b\in I^mV.
    \]
    Therefore
    \[
        \operatorname{ord}_I(a+b)
        \geq
        \min\{\operatorname{ord}_I(a),\operatorname{ord}_I(b)\}.
    \]
    Since \(0<\gamma<1\), this gives
    \[
        \|a+b\|_I
        \leq
        \max\{\|a\|_I,\|b\|_I\}.
    \]
    Thus
    \[
        d_I(v,u)
        =
        \|v-u\|_I
        =
        \|a+b\|_I
        \leq
        \max\{\|a\|_I,\|b\|_I\}
        =
        \max\{d_I(v,w),d_I(w,u)\}.
    \]
    Hence \(d_I\) is a non-archimedean metric on \(V\).
\end{proof}

\begin{proposition}[Cauchy sequences and convergence]
    Let \((v_i)_{i\geq 1}\) be a sequence of elements of \(V\).
    Then:
    \begin{enumerate}
        \item The sequence \((v_i)\) is Cauchy with respect to \(d_I\) if and only if
        for every integer \(m>0\), there exists an integer \(n>0\) such that
        \[
            v_i-v_j\in I^mV
        \]
        for all \(i,j>n\).

        \item The sequence \((v_i)\) converges to \(v\in V\) with respect to \(d_I\) if
        and only if for every integer \(m>0\), there exists an integer \(n>0\) such that
        \[
            v_i-v\in I^mV
        \]
        for all \(i>n\).
    \end{enumerate}
\end{proposition}

\begin{proof}
    We prove the first statement. The second is analogous.

    Suppose first that \((v_i)\) is Cauchy. Let \(m>0\). Since
    \[
        \gamma^m>0,
    \]
    there exists \(n>0\) such that
    \[
        d_I(v_i,v_j)<\gamma^m
    \]
    for all \(i,j>n\). This means
    \[
        \|v_i-v_j\|_I<\gamma^m.
    \]
    By the definition of the \(I\)-adic norm, this implies
    \[
        v_i-v_j\in I^mV.
    \]

    Conversely, suppose that for every \(m>0\), there exists \(n>0\) such that
    \[
        v_i-v_j\in I^mV
    \]
    for all \(i,j>n\). Let \(\varepsilon>0\). Since \(0<\gamma<1\), we may choose
    \(m>0\) such that
    \[
        \gamma^m<\varepsilon.
    \]
    By assumption, there exists \(n>0\) such that
    \[
        v_i-v_j\in I^mV
    \]
    for all \(i,j>n\). Hence
    \[
        d_I(v_i,v_j)=\|v_i-v_j\|_I\leq \gamma^m<\varepsilon.
    \]
    Therefore \((v_i)\) is a Cauchy sequence.
\end{proof}

\begin{lemma}
    With respect to the \(I\)-adic metrics, the following maps are continuous:
    \[
        V\times V \longrightarrow V,
        \qquad
        (a,b)\longmapsto a+b,
    \]
    \[
        R\times V \longrightarrow V,
        \qquad
        (r,v)\longmapsto rv,
    \]
    \[
        R\times R \longrightarrow R,
        \qquad
        (r,s)\longmapsto r+s,
    \]
    \[
        R\times R \longrightarrow R,
        \qquad
        (r,s)\longmapsto rs.
    \]
    Moreover, if \(W\) is another \(R\)-module endowed with its \(I\)-adic metric, then
    every \(R\)-module homomorphism
    \[
        f\in \operatorname{Hom}_R(V,W)
    \]
    is continuous.
\end{lemma}

\begin{proof}
    We indicate the main estimates.

    For addition in \(V\), one has
    \[
        \|(a+b)-(a'+b')\|_I
        =
        \|(a-a')+(b-b')\|_I
        \leq
        \max\{\|a-a'\|_I,\|b-b'\|_I\}.
    \]
    Hence addition is continuous.

    For scalar multiplication, suppose \(r-r'\in I^m\) and \(v-v'\in I^mV\).
    Then
    \[
        rv-r'v'
        =
        r(v-v')+(r-r')v'.
    \]
    Since
    \[
        r(v-v')\in I^mV,
        \qquad
        (r-r')v'\in I^mV,
    \]
    we get
    \[
        rv-r'v'\in I^mV.
    \]
    Thus scalar multiplication is continuous.

    The proofs for addition and multiplication in \(R\) are the special cases obtained
    by taking \(V=R\).

    Finally, let \(f:V\to W\) be an \(R\)-module homomorphism. If
    \[
        v-v'\in I^mV,
    \]
    then
    \[
        f(v)-f(v')=f(v-v')\in I^mW.
    \]
    Hence \(f\) is continuous.
\end{proof}

Before discussing completeness, one should first check that all the maps considered
are well-defined. Once this is done, the continuity statements from the previous 
show that the \(I\)-adic topology is compatible with the algebraic operations.

\begin{definition}[\(I\)-adic completeness]
    Let \(R\) be a ring and let \(I\) be an ideal of \(R\). Suppose that the
    \(I\)-adic metric \(d_I\) is defined on \(R\), namely
    \[
        \bigcap_{n\geq 1} I^n = 0.
    \]
    We say that \(R\) is complete with respect to \(I\) if the metric space
    \[
        (R,d_I)
    \]
    is complete.

    Similarly, if \(V\) is an \(R\)-module and the \(I\)-adic metric \(d_I\) is
    defined on \(V\), namely
    \[
        \bigcap_{n\geq 1} I^nV=0,
    \]
    then we say that \(V\) is complete with respect to \(I\) if the metric space
    \[
        (V,d_I)
    \]
    is complete.
\end{definition}

\begin{theorem}
    Suppose that the \(I\)-adic metric \(d_I\) is defined on every finitely generated
    \(R\)-module. Let \(V\) be a finitely generated \(R\)-module. Then every submodule
    of \(V\) is closed with respect to the \(I\)-adic metric on \(V\).
\end{theorem}

\begin{proof}
    Let \(W\leq V\) be a submodule. Consider the quotient map
    \[
        f:V\longrightarrow V/W.
    \]
    Since \(V\) is finitely generated, the quotient \(V/W\) is also finitely generated.
    By assumption, the \(I\)-adic metric is defined on \(V/W\).

    The quotient map \(f\) is an \(R\)-module homomorphism, and hence is continuous
    with respect to the \(I\)-adic metrics. Moreover,
    \[
        W=f^{-1}(0).
    \]
    Since \(\{0\}\) is closed in the metric space \(V/W\), its inverse image under the
    continuous map \(f\) is closed in \(V\). Therefore \(W\) is closed in \(V\).
\end{proof}

\begin{theorem}
    Let \(V\) be a finitely generated \(R\)-module. Assume that the \(I\)-adic
    metric is defined on \(V\). If \(R\) is complete with respect to the
    \(I\)-adic metric, then \(V\) is complete with respect to the \(I\)-adic metric.
\end{theorem}

\begin{proof}
    Let
    \[
        v_1,\dots,v_s
    \]
    be a finite set of generators of \(V\), and let
    \[
        (w_n)_{n\geq 1}
    \]
    be a Cauchy sequence in \(V\).

    Since \((w_n)\) is Cauchy, there exists a nondecreasing sequence of nonnegative
    integers
    \[
        m(n)\qquad (n\geq 1)
    \]
    such that
    \[
        m(n)\longrightarrow \infty
    \]
    and
    \[
        w_i-w_j\in I^{m(n)}V
    \]
    for all \(i,j>n\). In particular, if
    \[
        \Delta_n:=w_{n+1}-w_n,
    \]
    then
    \[
        \Delta_n\in I^{m(n)}V.
    \]

    Write
    \[
        w_1=\sum_{t=1}^s b_{1t}v_t,
        \qquad b_{1t}\in R.
    \]
    Since \(\Delta_n\in I^{m(n)}V\), and since \(v_1,\dots,v_s\) generate \(V\), we may
    choose elements
    \[
        a_{nt}\in I^{m(n)}
        \qquad (1\leq t\leq s)
    \]
    such that
    \[
        \Delta_n=w_{n+1}-w_n=\sum_{t=1}^s a_{nt}v_t.
    \]
    For \(n\geq 1\), define
    \[
        b_{nt}:=b_{1t}+\sum_{j=1}^{n-1}a_{jt}.
    \]
    We claim that
    \[
        w_n=\sum_{t=1}^s b_{nt}v_t
    \]
    for every \(n\geq 1\). Indeed, this is clear for \(n=1\), and if it holds for
    \(n\), then
    \[
        w_{n+1}
        =
        w_n+\Delta_n
        =
        \sum_{t=1}^s b_{nt}v_t+\sum_{t=1}^s a_{nt}v_t
        =
        \sum_{t=1}^s (b_{nt}+a_{nt})v_t
        =
        \sum_{t=1}^s b_{n+1,t}v_t.
    \]
    Hence the claim follows by induction.

    We now show that each sequence \((b_{nt})_{n\geq 1}\) is Cauchy in \(R\). Fix
    \(t\) and let \(\ell\geq 1\). Since \(m(n)\to\infty\), there exists \(N\) such that
    \[
        m(N)\geq \ell.
    \]
    If \(q>p\geq N\), then
    \[
        b_{qt}-b_{pt}
        =
        \sum_{j=p}^{q-1}a_{jt}.
    \]
    Since \(m(j)\geq m(p)\geq \ell\) for \(j\geq p\), we have
    \[
        a_{jt}\in I^{m(j)}\subseteq I^\ell.
    \]
    Therefore
    \[
        b_{qt}-b_{pt}\in I^\ell.
    \]
    Thus \((b_{nt})_{n\geq 1}\) is Cauchy in \(R\).

    Since \(R\) is complete, for each \(t=1,\dots,s\), there exists \(b_t\in R\) such that
    \[
        b_{nt}\longrightarrow b_t
    \]
    in the \(I\)-adic metric on \(R\). Define
    \[
        w:=\sum_{t=1}^s b_tv_t\in V.
    \]

    It remains to prove that \(w_n\to w\). Let \(\ell\geq 1\). For each \(t\), since
    \(b_{nt}\to b_t\), there exists \(N_t\) such that
    \[
        b_{nt}-b_t\in I^\ell
    \]
    for all \(n\geq N_t\). Put
    \[
        N=\max\{N_1,\dots,N_s\}.
    \]
    Then for \(n\geq N\),
    \[
        w_n-w
        =
        \sum_{t=1}^s (b_{nt}-b_t)v_t
        \in I^\ell V.
    \]
    Since \(\ell\geq 1\) was arbitrary, this shows that
    \[
        w_n\longrightarrow w
    \]
    in the \(I\)-adic metric on \(V\).

    Therefore every Cauchy sequence in \(V\) converges in \(V\). Hence \(V\) is complete.
\end{proof}
